\documentclass[11pt]{article}

\usepackage[T1]{fontenc}
\usepackage[margin=1.1in]{geometry}
\usepackage{amsmath,amssymb,amsthm}
\usepackage{enumitem}
\usepackage[colorlinks=true,linkcolor=blue,citecolor=blue,urlcolor=blue]{hyperref}

\newtheorem{theorem}{Theorem}[section]
\newtheorem{proposition}[theorem]{Proposition}
\newtheorem{lemma}[theorem]{Lemma}
\newtheorem{corollary}[theorem]{Corollary}
\newtheorem{definition}[theorem]{Definition}
\theoremstyle{remark}
\newtheorem{remark}[theorem]{Remark}
\newtheorem{example}[theorem]{Example}

\newcommand{\R}{\mathbb{R}}
\newcommand{\B}{\mathbb{B}}
\newcommand{\X}{X}
\newcommand{\conv}{\operatorname{conv}}
\newcommand{\cconv}[1]{\overline{\conv}\,#1}
\newcommand{\dist}{\operatorname{dist}}
\newcommand{\sca}[2]{\left\langle #1,\, #2\right\rangle}
\newcommand{\norm}[1]{\left\lVert #1\right\rVert}
\newcommand{\Mx}{M_{\X}}
\newcommand{\Cx}{C_{\X}}

\title{Notes on the proofs in\\ \emph{Sets Reconstructible with Medial Axis}\\[2mm]
\large prepared for a formalization in Lean~4 / mathlib\thanks{Notes drafted with the assistance
of Claude (Anthropic) during the planning stage of the formalization. All arguments below were
re-derived and checked by hand at the level of detail a proof assistant requires; remaining
errors are ours.}}
\author{}
\date{June 2026}

\begin{document}
\maketitle

\begin{abstract}
We are formalizing Theorem~6 of \cite{Bia} in the Lean~4 proof assistant, on top of the
\texttt{mathlib} library. Machine checking forces every implicit step of a proof to be made fully
explicit, and these notes record what that process surfaced. Concretely, we document:
(i)~one classical background fact which the paper uses through a citation
(the Motzkin--Fremlin inclusion $\Cx\subseteq\overline{\Mx}$) and for which we give a
self-contained elementary proof that avoids Brouwer's fixed point theorem (which is not
available in \texttt{mathlib});
(ii)~a small toolbox of lemmas (a ``ball membership along a ray'' criterion and an
``inflation dichotomy'') through which all the ball-inflation arguments of the paper factor;
(iii)~a reorganization in which Propositions~3 and~4 of \cite{Bia} merge into a single statement
whose proof does not use Motzkin's theorem on Chebyshev sets at all, and turns out to be
entirely free of limit arguments;
(iv)~a complete replacement for the proof of the forward direction of Theorem~6, where the
argument sketched in \cite{Bia} has a genuine gap.
None of this affects the validity of the \emph{statements} in \cite{Bia}: Theorem~6 is correct as
stated. The notes are written in ordinary mathematical language and require no knowledge of Lean.
As of this version, every proof below \emph{except} Appendix~\ref{app:fremlin} has been
machine-verified: Theorem~6 is fully checked in Lean modulo the single classical input
\eqref{eq:fremlin}.
\end{abstract}

\section{Setting, notation, and the result being formalized}

Throughout, $E$ is a finite-dimensional real inner product space (the formalization is carried
out at this level of generality; $E=\R^n$ is the motivating case), and $\X\subseteq E$ is a
nonempty closed set. Following \cite{Bia} we write
\[
 d(a) \;=\; d(a,\X)\;=\;\inf_{x\in \X}\norm{a-x},\qquad
 m(a)\;=\;\{x\in \X \mid \norm{a-x}=d(a)\},
\]
$\B(a,r)$ for the \emph{open} ball, and
\[
 \Mx=\{a\in E\mid \#m(a)>1\},\qquad
 \Cx=\{\text{centres of inclusion-maximal open balls contained in } E\setminus \X\}.
\]
A point $p\in E\setminus\X$ is \emph{reconstructible} if
$p\in\bigcup_{a\in \Mx}\B(a,d(a))$, and $E\setminus\X$ is reconstructible if all its points are.
A hyperplane $L$ is \emph{supporting} for a set $S$ if $S\cap L\neq\emptyset$ and $S$ lies in one
of the two closed half-spaces bounded by $L$; then $L^{+}$ denotes the open half-space disjoint
from $S$. We represent hyperplanes by a unit normal: $L=\{x\mid \sca{\eta}{x}=c\}$ with
$S\subseteq\{\sca{\eta}{\cdot}\le c\}$ and $L^{+}=\{\sca{\eta}{\cdot}>c\}$. We call a supporting
hyperplane $L$ of $\cconv \X$ \emph{defective} if $\X\cap L\neq \cconv \X\cap L$; equivalently
(since $\X\cap L\subseteq\cconv\X\cap L$ always) if there exists a \emph{defect witness}
\[
 w\in (\cconv \X\cap L)\setminus \X .
\]

The target is:

\begin{theorem}[{\cite[Theorem 6]{Bia}}]\label{thm:main}
Let $\X\subseteq E$ be closed and nonempty, and let $\mathcal L$ be the family of all defective
supporting hyperplanes of $\cconv \X$. Then $E\setminus\X$ is reconstructible if and only if
\[
E\setminus\X\;\subseteq\; \cconv \X\,\cup \bigcup_{L\in\mathcal L} L^{+}.
\]
\end{theorem}

(The paper says ``let $\mathcal L$ be \emph{a} family\dots''; for the forward implication
$\mathcal L$ must be the family of \emph{all} defective supporting hyperplanes, and we formalize
the statement in that form.)

\subsection*{Summary of findings}

\begin{center}
\renewcommand{\arraystretch}{1.25}
\begin{tabular}{p{0.27\textwidth}p{0.66\textwidth}}
\hline
Result in \cite{Bia} & Status in the formalization \\
\hline
Lemma 1 & Proved; the step ``$c$ can be approximated by points of $\Mx$'' is the
Motzkin--Fremlin theorem $\Cx\subseteq\overline{\Mx}$, which we must prove ourselves
(\S\ref{sec:fremlin}, Appendix~\ref{app:fremlin}); in the first stage of the formalization this
single classical fact is assumed.\\
Proposition 2 & Proved; one definition repaired (Remark~\ref{rem:prop2sup}) and Carath\'eodory's
theorem is not needed (Remark~\ref{rem:prop2cara}).\\
Propositions 3, 4 & Merged into one statement (Theorem~\ref{thm:unified}) with a shorter proof:
Motzkin's theorem on Chebyshev sets, the cone at infinity, and the two-nearest-points analysis
are no longer used; moreover the proof is entirely limit-free — a single explicit value of the
lift parameter suffices (Remark~\ref{rem:swallow}).\\
Corollary 5 & As written it invokes the \emph{proof} (not the statement) of Proposition~4 and
does not literally follow from it; we give a direct ``one-shot'' proof needing no limits
(\S\ref{sec:cor5}).\\
Theorem 6, ``$\Leftarrow$'' & Assembly of the above. \\
Theorem 6, ``$\Rightarrow$'' & The sketch in \cite{Bia} has a gap (the segment
$\left]p,q\right[$ may miss $\cconv\X$ entirely, and a separating hyperplane through the entry
point need not exist); we give a complete replacement proof via the metric projection onto
$\cconv \X$ (\S\ref{sec:forward}).\\
\hline
\end{tabular}
\end{center}

\section{Toolbox}\label{sec:toolbox}

All inflation arguments in \cite{Bia} grow a one-parameter family of balls that share a tangency
point. The following trivial identity is the workhorse that replaces every limit computation
``$p$ eventually falls into the growing balls'' by an explicit inequality. (In a formalization
such a reduction matters enormously; but we believe it also makes the paper proofs shorter.)

\begin{lemma}[ball membership along a ray]\label{lem:swallow}
Let $x,p\in E$, let $v$ be a unit vector and $\rho>0$. Then
\[
 p\in\B(x+\rho v,\rho)\iff 2\rho\,\sca{v}{p-x}>\norm{p-x}^2 .
\]
\end{lemma}

\begin{proof}
Expand $\norm{p-x-\rho v}^2<\rho^2$.
\end{proof}

\begin{lemma}[tangent families]\label{lem:tangent}
Fix $x\in E$ and a unit vector $v$, and write $\B_t:=\B(x+tv,t)$ for $t>0$ (the family of open
balls tangent to the hyperplane through $x$ with normal $v$, at the point $x$). Then
\begin{enumerate}[label=(\alph*),itemsep=1pt]
\item $0<s\le t\implies \B_s\subseteq \B_t$;
\item $\bigcup_{0<t<T}\B_t=\B_T$ for every $T>0$;
\item $\bigcup_{t>0}\B_t=\{y\mid\sca{v}{y-x}>0\}$, the open half-space.
\end{enumerate}
\end{lemma}

\begin{proof}
All three are immediate from Lemma~\ref{lem:swallow}: membership of $y$ in $\B_t$ reads
$2t\sca{v}{y-x}>\norm{y-x}^2$, in which the left side is nonnegative-increasing in $t$ once
$\sca{v}{y-x}>0$, and which for $y\ne x$ can hold only if $\sca{v}{y-x}>0$.
\end{proof}

We record the elementary facts about maximal balls; recall a ball
$\B(c,r)\subseteq E\setminus \X$, $r>0$, is \emph{maximal} if
$\B(c,r)\subseteq\B(c',r')\subseteq E\setminus\X$ (with $r'>0$) forces $c'=c$, $r'=r$.

\begin{lemma}\label{lem:maxballs}
Let $\X$ be closed and nonempty.
\begin{enumerate}[label=(\alph*),itemsep=1pt]
\item For balls in a normed space, $\B(c,r)\subseteq\B(c',r')$ iff $\norm{c-c'}\le r'-r$.
\item A maximal ball centred at $c$ has radius $d(c)$; thus
$\Cx=\{c\mid \B(c,d(c))\text{ is maximal}\}$.
\item $\Mx\subseteq \Cx$.
\item If $a\in \X$ then $m(a)=\{a\}$; hence $\Mx\cap \X=\emptyset$ and $d(a)>0$ for
$a\in\Mx\cup\Cx$.
\end{enumerate}
\end{lemma}

\begin{proof}[Proof sketch]
(b) $\B(c,r)\subseteq E\setminus\X$ iff $r\le d(c)$. (c) Let $a\in\Mx$ with distinct
$x,y\in m(a)$ and suppose $\B(a,d(a))\subseteq\B(c',r')\subseteq E\setminus\X$. Since
$x\in \X$, $x\notin\B(c',r')$, so $\norm{x-c'}\ge r'$; on the other hand by (a),
$\norm{x-c'}\le\norm{x-a}+\norm{a-c'}\le d(a)+(r'-d(a))=r'$. So equality holds in the triangle
inequality, which in a Euclidean (strictly convex) space forces $c'$ to lie on the ray from $x$
through $a$, at distance $r'-d(a)$ beyond $a$; the same holds with $y$ in place of $x$. If
$r'>d(a)$ the two descriptions of $c'$ give $(a-x)/d(a)=(a-y)/d(a)$, i.e.\ $x=y$, a
contradiction; hence $r'=d(a)$ and $c'=a$.
\end{proof}

The next lemma is the common engine behind Propositions 2--4 and Corollary 5 of \cite{Bia}.
It is the precise form of ``inflate a ball keeping the tangency point until it either becomes
maximal or sweeps out a half-space''.

\begin{lemma}[inflation dichotomy]\label{lem:dicho}
Let $\X$ be closed and nonempty, $x_0\in \X$, $v$ a unit vector, and suppose
$\B(x_0+t_0v,t_0)\cap \X=\emptyset$ for some $t_0>0$. Put
\[
 T:=\sup\{t>0\mid \B(x_0+tv,t)\cap \X=\emptyset\}\;\in\;[t_0,\infty].
\]
\begin{enumerate}[label=(\roman*),itemsep=1pt]
\item If $T=\infty$, then $\X\cap\{y\mid\sca{v}{y-x_0}>0\}=\emptyset$.
\item If $T<\infty$, then $c:=x_0+Tv\in\Cx$, with $d(c)=T$ and $x_0\in m(c)$.
\item If moreover $p\notin \X$, $x_0\in m(p)$ and $v=(p-x_0)/d(p)$, then $t_0=d(p)$ qualifies,
and in case~(ii) one has $p\in\B(c,d(c))$.
\end{enumerate}
\end{lemma}

\begin{proof}
By Lemma~\ref{lem:tangent}(a) the set of admissible $t$ is an interval, so every $t<T$ is
admissible. (i)~is Lemma~\ref{lem:tangent}(c). For (ii): by Lemma~\ref{lem:tangent}(b),
$\B(c,T)=\bigcup_{t<T}\B_t$ is disjoint from $\X$, so $d(c)\ge T$; since $x_0\in \X$ and
$\norm{c-x_0}=T$ we get $d(c)=T$ and $x_0\in m(c)$. Maximality: if
$\B(c,T)\subseteq\B(c',r')\subseteq E\setminus\X$ then, exactly as in the proof of
Lemma~\ref{lem:maxballs}(c), $\norm{x_0-c'}=r'$ with equality in the triangle inequality through
$c$, so $c'=x_0+r'v$; if $r'>T$ this contradicts the definition of $T$, hence $r'=T$, $c'=c$.
For (iii): $\B(x_0+d(p)v,d(p))=\B(p,d(p))$ is disjoint from $\X$, so $T\ge d(p)$; and by
Lemma~\ref{lem:swallow}, $p\in\B(c,T)$ iff $2T\sca{v}{p-x_0}=2T\,d(p)>d(p)^2$, which holds
because $T\ge d(p)>\tfrac12 d(p)$.
\end{proof}

\section{Lemma 1 and the Motzkin--Fremlin theorem}\label{sec:fremlin}

The proof of \cite[Lemma~1]{Bia} uses that every $c\in\Cx$ ``can be approximated by a sequence of
points $x_n\in\Mx$'', i.e.
\begin{equation}\label{eq:fremlin}
 \Mx\;\subseteq\;\Cx\;\subseteq\;\overline{\Mx},
\end{equation}
which the paper quotes from the literature (Motzkin; Fremlin \cite{Fre}). For the formalization
this inclusion must itself be proved, and it sits on the critical path of \emph{everything}: all
the reconstruction arguments produce centres of maximal balls, and a maximal ball may touch $\X$
at a single point, so the passage from $\Cx$ back to $\Mx$ is genuinely needed and genuinely
nontrivial. The textbook proofs use Brouwer's fixed point theorem or degree theory, neither of
which is currently available in \texttt{mathlib}. In Appendix~\ref{app:fremlin} we give a
self-contained proof of $\Cx\subseteq\overline{\Mx}$ which uses nothing beyond compactness and an
explicit Euler-polygon construction; we believe it may be of independent interest. (In the first
stage of the formalization, \eqref{eq:fremlin} is assumed as a clearly marked placeholder and
everything else is proved unconditionally; in the second stage we formalize
Appendix~\ref{app:fremlin}.)

Given \eqref{eq:fremlin}, Lemma~1 of \cite{Bia} is straightforward, exactly as in the paper:

\begin{lemma}[{\cite[Lemma~1]{Bia}}]\label{lem:one}
Let $p\in E\setminus \X$. Then $p$ is reconstructible if and only if there is $c\in\Cx$ with
$p\in\B(c,d(c))$.
\end{lemma}

\begin{proof}
($\Rightarrow$) $\Mx\subseteq\Cx$ by Lemma~\ref{lem:maxballs}(c).
($\Leftarrow$) By \eqref{eq:fremlin} pick $a_n\in\Mx$, $a_n\to c$. Since $d$ is $1$-Lipschitz,
$d(a_n)\to d(c)$ and $\norm{p-a_n}\to\norm{p-c}<d(c)$, so $\norm{p-a_n}<d(a_n)$ for large $n$.
\end{proof}

\section{Proposition 2}\label{sec:prop2}

\begin{remark}[a repaired definition]\label{rem:prop2sup}
The proof of \cite[Proposition~2]{Bia} introduces, for $p_t:=t(p-x_p)+x_p$,
\[
 r:=\sup\{t>0\mid \B(p_t,d(p_t))\cap \X=\emptyset\},
\]
and argues by cases on $r<\infty$. As literally written the defining condition is vacuous:
$\B(a,d(a))\cap \X=\emptyset$ holds for \emph{every} $a$, so $r=\infty$ always. What is meant
(and what the surrounding text uses, e.g.\ ``it is the union of the family of balls'') is the
family of balls \emph{tangent to $\X$ at $x_p$}, i.e.\ $\B(p_t,t\,d(p))$, equivalently the
condition $d(p_t)=t\,d(p)$. Lemma~\ref{lem:dicho} is exactly this corrected statement, packaged
once and for all; the same correction applies wherever the ``directional reaching radius''
$r_v(x):=\sup\{\rho\ge0\mid\B(\rho v+x,\,\rho)\cap \X=\emptyset\}$ of \cite{Bia} is used, which
matches our $T$ in Lemma~\ref{lem:dicho}.
\end{remark}

\begin{proposition}[{\cite[Proposition~2]{Bia}}]\label{prop:two}
If $p\in\conv \X\setminus \X$, then $p$ is reconstructible.
\end{proposition}

\begin{proof}
Pick $x_p\in m(p)$ (nonempty: $\X$ closed nonempty, $E$ finite-dimensional) and let
$v=(p-x_p)/d(p)$. Apply Lemma~\ref{lem:dicho}(iii). If $T<\infty$ we obtain $c\in\Cx$ with
$p\in\B(c,d(c))$, and Lemma~\ref{lem:one} concludes. If $T=\infty$, then
$\X\subseteq K:=\{y\mid \sca{v}{y-x_p}\le0\}$. But $K$ is convex, so $\conv \X\subseteq K$, and
$p\in\conv \X$ gives $\sca{v}{p-x_p}\le0$; since $\sca{v}{p-x_p}=d(p)>0$, this is a
contradiction.
\end{proof}

\begin{remark}\label{rem:prop2cara}
The paper's proof invokes Carath\'eodory's lemma to write $p$ as a convex combination of points
of $\X$ and then computes. This is not needed: the half-space $K$ is convex and contains $\X$, so
it contains $\conv \X\ni p$ --- the computation with the convex combination in \cite{Bia} is
precisely the proof of this containment in the special case. (In the formalization this saves a
dependency on the Carath\'eodory machinery.)
\end{remark}

\section{A unified Proposition 3/4}\label{sec:unified}

This section contains the main reorganization. Recall \cite[Proposition~3]{Bia}: if $L$ is a
supporting hyperplane of $\conv \X$ with $\X\cap L$ \emph{non-convex}, then every $p\in L^{+}$
is reconstructible; and \cite[Proposition~4]{Bia}: the same conclusion when $L$ is a supporting
hyperplane of $\cconv \X$ with $\X\cap L\neq\cconv \X\cap L$. The proof of Proposition~3 in
\cite{Bia} starts from Motzkin's theorem (a closed non-convex set has a nonempty medial axis,
applied to $\X\cap L$ inside $L$) to obtain a point $q$ with two nearest points $a\neq b$ in
$\X\cap L$, and the proof of Proposition~4 reduces its convex case to a separate ladder
construction.

Two observations let us merge the two propositions and shorten the proof considerably.

\begin{enumerate}[label=(\Alph*),itemsep=2pt]
\item The hypothesis of Proposition 3 implies that of Proposition 4: if $\X\cap L$ is
non-convex, then $\X\cap L\subsetneq\conv(\X\cap L)\subseteq\cconv\X\cap L$, so $L$ is
defective. Hence a single statement with the ``defective'' hypothesis covers both.
\item In the ladder construction, the role played in \cite{Bia} by the pair $a,b$ of nearest
points (Proposition 3) or by approximating convex combinations $a_\nu,b_\nu$ (Proposition 4) can
be played directly by a \emph{defect witness} $w\in(\cconv\X\cap L)\setminus \X$: in the
``infinite radius'' case the half-space produced by Lemma~\ref{lem:dicho}(i) is closed and
convex, hence contains $\cconv \X\ni w$, and evaluating the half-space inequality at $w$ gives an
immediate contradiction. \emph{Consequently Motzkin's theorem is not needed anywhere in the proof
of Theorem~6.} (It remains an interesting target in its own right; see
Remark~\ref{rem:motzkin}.)
\end{enumerate}

\begin{theorem}[unified Propositions 3 and 4]\label{thm:unified}
Let $\X$ be closed and nonempty and let $L=\{x\mid\sca{\eta}{x}=c\}$, $\norm{\eta}=1$, satisfy
\[
\X\subseteq\{x\mid \sca{\eta}{x}\le c\},\qquad
\exists\, w\in \cconv\X\cap L,\ w\notin \X .
\]
Then every $p$ with $\sca{\eta}{p}>c$ is reconstructible.
\end{theorem}

\begin{proof}
Translating, we may assume $c=0$; write $\gamma:=\sca{\eta}{p}>0$, and fix \emph{once and for
all} the lift parameter
\[
 t\;:=\;\frac{\norm{p-w}^2}{2\gamma}+1,\qquad\text{so that}\qquad
 \norm{p-w}^2\;=\;2t\gamma-2\gamma .
\]
Consider the lifted point $w_t:=w+t\eta$; since $\sca{\eta}{w_t}=t>0$ while
$\X\subseteq\{\sca{\eta}{\cdot}\le0\}$, we have $w_t\notin \X$, so
\[
 x_t\in m(w_t),\qquad D:=d(w_t)=\norm{w_t-x_t}>0,\qquad v:=\frac{w_t-x_t}{D}
\]
are well defined. Abbreviate
\[
 S:=w-x_t,\qquad P:=p-w,\qquad \beta:=-\sca{\eta}{x_t}\;\ge\;0
\]
(the inequality because $x_t\in \X$). Everything rests on two instances of the same expansion,
using $w_t-x_t=t\eta+S$, $\sca{\eta}{S}=\beta$ and $\sca{\eta}{P}=\gamma$:
\begin{align}
 \sca{w_t-x_t}{\,w-x_t}&=\sca{t\eta+S}{S}\;=\;t\beta+\norm S^2,\label{eq:expandw}\\
 N\;:=\;\sca{w_t-x_t}{\,p-x_t}&=\sca{t\eta+S}{P+S}\;=\;t\gamma+t\beta+\sca{S}{P}+\norm S^2 .
 \label{eq:expandp}
\end{align}
The ball $\B(w_t,D)$ is the tangent ball of radius $D$ at $x_t$ in the direction $v$ and avoids
$\X$; apply Lemma~\ref{lem:dicho}.

\emph{The case $T=\infty$ is impossible.} By Lemma~\ref{lem:dicho}(i),
$\X\subseteq K:=\{y\mid\sca{v}{y-x_t}\le0\}$; $K$ is closed and convex, so
$\cconv \X\subseteq K$, and in particular $\sca{w_t-x_t}{w-x_t}\le0$. By \eqref{eq:expandw}
this forces $t\beta+\norm S^2\le0$, hence $S=0$, i.e.\ $x_t=w$ --- contradicting $x_t\in \X$,
$w\notin \X$.

\emph{The case $T<\infty$, and the capture of $p$.} By Lemma~\ref{lem:dicho}(ii),
$c_T:=x_t+Tv\in\Cx$ with $d(c_T)=T\ge D$. By Lemmas~\ref{lem:swallow} and~\ref{lem:one} it
remains to check
\begin{equation}\label{eq:goal}
 \norm{p-x_t}^2\;<\;2T\sca{v}{p-x_t}\;=\;\frac{2T}{D}\,N .
\end{equation}
First, $N\ge0$: by Cauchy--Schwarz $\sca{S}{P}\ge-\norm S\norm P$, by
$(\norm S-\tfrac12\norm P)^2\ge0$ we have $\norm S^2-\norm S\norm P\ge-\tfrac14\norm P^2$, and
$\norm P^2\le 2t\gamma$ by the choice of $t$; hence
\[
 N\;\ge\;t\gamma+t\beta-\tfrac14\norm P^2\;\ge\;t\gamma-\tfrac12\,t\gamma\;>\;0 .
\]
Consequently $\tfrac{2T}{D}N\ge2N$, and it suffices to compare $2N$ with
$\norm{p-x_t}^2=\norm{P+S}^2=\norm P^2+2\sca{S}{P}+\norm S^2$. The cross terms cancel
\emph{exactly} against \eqref{eq:expandp}:
\[
 2N-\norm{p-x_t}^2\;=\;2t\gamma+2t\beta+\norm S^2-\norm P^2
 \;\ge\;2t\gamma-\norm P^2\;=\;2\gamma\;>\;0 .
\]
This proves \eqref{eq:goal}; so $p\in\B(c_T,d(c_T))$ with $c_T\in\Cx$, and Lemma~\ref{lem:one}
finishes the proof.
\end{proof}

\begin{remark}[on the final limit step in \cite{Bia}]\label{rem:swallow}
The proofs of Propositions~3 and~4 in \cite{Bia} conclude with: \emph{``For any $p$ in $L^{+}$
and a sequence of points $c_n\in\Cx$ with $\norm{c_n}\to\infty$, $c_n/\norm{c_n}\to\eta$ we will
have $d(c_n)>\norm{p-c_n}$ for $n$ large enough.''} The two stated properties of the sequence
are not by themselves sufficient for the conclusion: writing $c_n=s_nu_n$ with $s_n\to\infty$,
$u_n\to\eta$, the available lower bound $d(c_n)\ge\sca{\eta}{c_n}$ beats
$\norm{p-c_n}\approx s_n-\sca{p}{u_n}$ only if $s_n(1-\sca{u_n}{\eta})$ stays small, which the
hypotheses do not guarantee for an \emph{arbitrary} such sequence. The sequences actually
\emph{constructed} in the paper do satisfy the conclusion, but proving it requires retaining the
tangency data of the maximal balls (tangency point and radius), which is what the proof above
does via Lemma~\ref{lem:swallow}. This is a typical example of a step that is morally right but
that a proof assistant (correctly) refuses without the extra data. Strikingly, once that data is
kept \emph{exactly}, all asymptotics evaporate: the cross terms $\sca{S}{P}$ cancel between the
two sides of \eqref{eq:goal}, a single explicit value of $t$ suffices, and neither the limit
$t\to\infty$ nor the coordinate estimates for $x_t$ in the proof of \cite[Proposition~4]{Bia}
(nor even the nearest point of $w$ itself) are needed. We found this simplification only while
machine-checking the proof.
\end{remark}

\begin{remark}[status of Proposition 3 and Motzkin's theorem]\label{rem:motzkin}
By observation (A) above, \cite[Proposition~3]{Bia} is a special case of
Theorem~\ref{thm:unified}; the formalization therefore does not state it separately, and
Motzkin's theorem (closed Chebyshev subsets of $\R^n$ are convex) leaves the dependency graph of
Theorem~6 entirely. We note that once $\Cx\subseteq\overline{\Mx}$ is available
(Appendix~\ref{app:fremlin}), Motzkin's theorem falls out for free, by the argument of
Proposition~\ref{prop:two}: if $\Mx=\emptyset$ then $\Cx=\emptyset$, so for any
$p\in\conv\X\setminus\X$ the dichotomy (Lemma~\ref{lem:dicho}(iii)) cannot terminate at a maximal
ball, hence produces a half-space through a nearest point separating $p$ from $\X$ ---
contradicting $p\in\conv \X$ as in Proposition~\ref{prop:two}. Hence a non-convex closed set has
points with at least two nearest points. We plan to include this as a corollary in the
formalization since it is of independent interest.
\end{remark}

\section{Corollary 5: a one-shot proof}\label{sec:cor5}

\begin{remark}
\cite[Corollary~5]{Bia} (every $p\in\cconv\X\setminus \X$ is reconstructible) is justified by
``any supporting hyperplane $L$ of $\cconv \X$ with $p\in L$ will raise a sequence of balls as
in the proof of Proposition~4''. Two points need attention. First, this invokes the
\emph{construction} inside the proof of Proposition~4 rather than its statement; moreover the
balls of that construction capture the open half-space $L^{+}$, while $p$ lies \emph{on} $L$, so
the capture of $p$ itself needs a separate argument. Second, the existence of a supporting
hyperplane through a boundary point requires a (standard but nonempty) argument when
$\cconv \X$ has empty interior. Both points are easy to settle, and in fact no limit
$t\to\infty$ is needed at all:
\end{remark}

\begin{corollary}[{\cite[Corollary~5]{Bia}}]\label{cor:five}
Every $p\in\cconv \X\setminus \X$ is reconstructible.
\end{corollary}

\begin{proof}
If $p\in\conv \X$, apply Proposition~\ref{prop:two}. So assume $p\in\cconv\X\setminus\conv\X$.

\emph{A supporting hyperplane through $p$.} If $\operatorname{int}\conv \X\neq\emptyset$ then
$\operatorname{int}\cconv\X=\operatorname{int}\conv \X\subseteq\conv\X$, so $p$ lies on the
boundary of the closed convex set $\cconv \X$ and a supporting hyperplane
$L=\{\sca{\eta}{\cdot}=c\}$ through $p$ exists by standard separation. If
$\operatorname{int}\conv \X=\emptyset$, then $\cconv \X$ lies in a proper affine subspace $A$;
any unit $\eta$ orthogonal to the direction of $A$ gives a hyperplane through $p$ containing
$\cconv\X$, which is in particular supporting. In both cases, after translation $c=0$,
\[
 \X\subseteq\{\sca{\eta}{\cdot}\le0\},\qquad \sca{\eta}{p}=0,\qquad p\in\cconv\X,\qquad
 p\notin \X .
\]

\emph{One-shot capture.} Fix any $t>0$ (say $t=1$) and run the construction of
Theorem~\ref{thm:unified} with defect witness $w:=p$ \emph{and target point $p$ itself}: let
$x_t\in m(p+t\eta)$, $D=d(p+t\eta)$, $v=(p+t\eta-x_t)/D$, and $S:=p-x_t\neq0$. The case
$T=\infty$ is impossible exactly as before (it would force $x_t=w=p\notin \X$). If $T<\infty$
we get $c_T\in\Cx$, $d(c_T)=T\ge D$, and the membership criterion for $p$ reads
\[
 2T\sca{v}{p-x_t}\;\ge\;\frac{2T}{D}\Bigl(t\sca{\eta}{S}+\norm S^2\Bigr)
 \;\ge\;2\norm S^2\;>\;\norm S^2=\norm{p-x_t}^2,
\]
using $\sca{\eta}{S}=-\sca{\eta}{x_t}\ge0$ and $T\ge D$. So $p\in\B(c_T,d(c_T))$, and
Lemma~\ref{lem:one} concludes. (No limit in $t$ was required; in the Lean development this case
and Theorem~\ref{thm:unified} are two instances of one lemma, with the hypotheses
$\sca{\eta}{p}\ge c$ and $\norm{p-w}^2\le 2t(\sca{\eta}{p}-c)$.)
\end{proof}

\section{Theorem 6}\label{sec:thm6}

\subsection{The backward implication}

Suppose $E\setminus \X\subseteq\cconv\X\cup\bigcup_{L\in\mathcal L}L^{+}$. Every
$p\in E\setminus\X$ then lies either in $\cconv \X$ (Corollary~\ref{cor:five}) or in $L^{+}$ for
some defective supporting hyperplane $L$ (Theorem~\ref{thm:unified}); in both cases $p$ is
reconstructible. \qed

\subsection{The forward implication: a gap and its repair}\label{sec:forward}

The argument sketched in \cite{Bia} is: given a reconstructible $p\notin\conv \X$ with
$p\in\B(a,d(a))$, $a\in\Mx$, pick $q\in m(a)$, connect $p$ and $q$ inside the ball, let
$[\bar a,\bar b[\,\subseteq\,]p,q[$ be the part of the segment inside $\cconv \X$, observe
$\bar a\in(\cconv\X\cap L)\setminus \X$ for any supporting $L\ni\bar a$, and finally ``pick $L$
to be a hyperplane that separates $p$ and $\cconv\X$''.

Two distinct difficulties arise when one tries to make this precise.

\begin{example}[the interval may not exist]\label{ex:twopoints}
Let $\X=\{(0,1),(0,-1)\}\subseteq\R^2$. Then $\Mx$ is the horizontal axis; take $a=(1,0)$,
$d(a)=\sqrt2$, and $p=(2,0)\in\B(a,\sqrt2)$, $p\notin\cconv\X=\{0\}\times[-1,1]$. For either
choice of $q\in m(a)=\{(0,\pm1)\}$ the open segment $]p,q[$ meets the vertical axis only in the
limit $q$ itself, so $]p,q[\,\cap\,\cconv\X=\emptyset$: the interval $[\bar a,\bar b[$ of the
proof is empty. (The theorem itself is fine here: $L=\{x_1=0\}$ is defective --- its hull section
is the segment, its $\X$-section two points --- and $p\in L^{+}$.)
\end{example}

\begin{example}[grazing]\label{ex:grazing}
Even when the entry point $\bar a$ exists, there need not be \emph{any} supporting hyperplane
through $\bar a$ having $p$ strictly on its open side. If the boundary of $\cconv\X$ near
$\bar a$ looks like the parabola $y=x^2$ near the origin and $p=(1,0)$, the unique supporting
line at $\bar a=(0,0)$ is $\{y=0\}$, which contains $p$ rather than separating it. So
``a hyperplane through $\bar a$ separating $p$ and $\cconv \X$'' may simply not exist, and the
two requirements on $L$ (defect witnessed at $\bar a\in L$; $p\in L^{+}$) cannot in general be
met by one hyperplane chosen this way.
\end{example}

The repair below produces the defect witness and the strict separation \emph{from the same
metric-projection construction}, using the second nearest point of $a$ (which the sketch in
\cite{Bia} never uses --- note that Example~\ref{ex:twopoints} shows one tangency point is not
enough\,!). Recall that in a finite-dimensional space the metric projection
$\pi:E\to\cconv \X$ onto a closed convex set is well defined (existence and uniqueness of the
nearest point), and is characterized by
\begin{equation}\label{eq:projchar}
 z=\pi(b)\iff z\in\cconv\X\ \text{ and }\ \forall x\in\cconv \X:\ \sca{b-z}{x-z}\le0 .
\end{equation}

\begin{theorem}[forward implication of Theorem~\ref{thm:main}]
Let $\X$ be closed and nonempty and let $p\in E\setminus\X$ be reconstructible with
$p\notin\cconv \X$. Then there is a defective supporting hyperplane $L$ of $\cconv\X$ with
$p\in L^{+}$.
\end{theorem}

\begin{proof}
Pick $a\in\Mx$ with $p\in\B:=\B(a,d(a))$ and \emph{two distinct} points $q,q'\in m(a)$, and set
\[
 q'':=\tfrac12(q+q') .
\]
By the parallelogram identity,
$\norm{a-q''}^2=d(a)^2-\tfrac14\norm{q-q'}^2<d(a)^2$, so $q''$ lies in the \emph{open} ball
$\B\subseteq E\setminus \X$; in particular $q''\notin \X$, while
$q''\in\conv\X\subseteq\cconv\X$.

Both endpoints of the segment $[p,q'']$ lie in the open convex set $\B$, hence so does the whole
segment. Since $q''\in\cconv \X$ and $p\notin\cconv\X$, and $\cconv\X$ is closed, the segment
has a \emph{first} point of $\cconv\X$ as seen from $p$: there is
$\bar a\in[p,q'']\cap\cconv \X$ with
\[
 [p,\bar a[\;\cap\;\cconv \X=\emptyset,\qquad \bar a\neq p .
\]
Note $\bar a\in\B$, so $\bar a\notin \X$. Let $u:=(p-\bar a)/\norm{p-\bar a}$.

For $b\in[p,\bar a[$ write $z_b:=\pi(b)$ and $g_b:=\norm{b-z_b}=\dist(b,\cconv\X)>0$. Since
$\bar a\in\cconv \X$,
\[
 \norm{z_b-\bar a}\le\norm{z_b-b}+\norm{b-\bar a}\le 2\norm{b-\bar a}\xrightarrow[b\to\bar a]{}0 .
\]
As $\bar a\notin \X$ and $\X$ is closed, we may fix $b\in[p,\bar a[$ with $z_b\notin \X$.

Now define the hyperplane by the unit normal $v:=(b-z_b)/g_b$ and level $c:=\sca{v}{z_b}$, i.e.\
$L:=\{x\mid \sca{v}{x}=c\}$. We check the three required properties.

\emph{$L$ supports $\cconv \X$.} By \eqref{eq:projchar}, $\sca{v}{x-z_b}\le0$ for all
$x\in\cconv \X$, i.e.\ $\cconv\X\subseteq\{\sca{v}{\cdot}\le c\}$; and $z_b\in\cconv\X\cap L$.
(A fortiori $\X\subseteq\{\sca{v}{\cdot}\le c\}$.)

\emph{$L$ is defective.} $z_b\in(\cconv \X\cap L)\setminus \X$ by the choice of $b$.

\emph{$p\in L^{+}$.} Since $b$ lies on the segment from $\bar a$ to $p$ we have
$b-\bar a=\norm{b-\bar a}\,u$ and $p-b=\norm{p-b}\,u$. Using $\sca{v}{\bar a}\le c$
(as $\bar a\in\cconv\X$),
\[
 \sca{v}{u}=\frac{\sca{v}{b-\bar a}}{\norm{b-\bar a}}
 =\frac{\sca{v}{b-z_b}+\sca{v}{z_b-\bar a}}{\norm{b-\bar a}}
 \;\ge\;\frac{g_b+0}{\norm{b-\bar a}}\;>\;0 ,
\]
hence
\[
 \sca{v}{p}-c=\sca{v}{p-z_b}=\sca{v}{p-b}+\sca{v}{b-z_b}
 =\norm{p-b}\,\sca{v}{u}+g_b\;>\;0 . \qedhere
\]
\end{proof}

\begin{remark}
The proof uses only the metric projection onto a closed convex set and no separation theorems;
it is insensitive to whether $\cconv\X$ has empty interior. In Example~\ref{ex:twopoints} it
produces $q''=(0,0)$, $\bar a=(0,0)$, $z_b=(0,0)$ and the defective vertical line, as expected.
Combining both implications with Corollary~\ref{cor:five} and Theorem~\ref{thm:unified} yields
Theorem~\ref{thm:main}.
\end{remark}

\section{Remarks on the formalization}\label{sec:lean}

The development has been carried out in Lean~4 over \texttt{mathlib}, for an arbitrary
finite-dimensional real inner product space; the definitions ($m$, $\Mx$, $\Cx$,
reconstructibility) are made for general metric spaces, with the geometric results at the
inner-product level. The structure mirrors these notes: the toolbox of \S\ref{sec:toolbox},
then Lemma~\ref{lem:one}, Proposition~\ref{prop:two}, Theorem~\ref{thm:unified},
Corollary~\ref{cor:five}, and the two implications of Theorem~\ref{thm:main}. As of this
version, \emph{all of the above is machine-checked} ($\approx1100$ lines in three files;
the main statement is \texttt{Metric.reconstructible\_iff}). The single classical input
\eqref{eq:fremlin} ($\Cx\subseteq\overline{\Mx}$) is assumed as an explicitly marked
placeholder, and Lean's axiom tracker confirms it is the \emph{only} unproved dependency of
Theorem~\ref{thm:main}; formalizing Appendix~\ref{app:fremlin} to discharge it is the next
step. The forward implication of Theorem~\ref{thm:main} (\S\ref{sec:forward}) is checked with
no assumptions at all. We will gladly share the Lean sources; comments on the mathematics in
these notes, and in particular on Appendix~\ref{app:fremlin}, are very welcome.

\appendix

\section{A Brouwer-free proof that \texorpdfstring{$\Cx\subseteq\overline{\Mx}$}{CX is contained
in the closure of MX}}\label{app:fremlin}

Throughout, $\X\subseteq E$ is closed and nonempty, $E$ finite-dimensional. We prove: \emph{the
centre of any maximal ball lies in the closure of $\Mx$.} The strategy: if $c\in\Cx$ had a
neighbourhood free of $\Mx$, the nearest-point map would be single-valued, hence continuous,
near $c$; we then inflate the maximal ball by following an explicit Euler polygon that always
moves directly away from the current nearest point, gaining distance from $\X$ at unit rate up
to an error controlled by the (uniform) continuity of the nearest-point map; a compactness
argument upgrades the almost-unit rate to unit rate, producing a ball that strictly contains the
maximal ball --- a contradiction.

\begin{lemma}\label{lem:footmap}
Let $K\subseteq E\setminus \X$ be compact with $K\cap\Mx=\emptyset$. Then every $a\in K$ has a
unique nearest point $P(a)$, and $P$ is uniformly continuous on $K$.
\end{lemma}

\begin{proof}
$m(a)\neq\emptyset$ by closedness of $\X$ and finite dimension; single-valuedness on $K$ is the
assumption. The graph of $m$ is closed (if $a_n\to a$, $x_n\in m(a_n)$, $x_n\to x$, then
$x\in \X$ and $\norm{a-x}=\lim d(a_n)=d(a)$ by continuity of $d$), and $m$ is locally bounded
($\norm{x}\le\norm a+d(a)$ for $x\in m(a)$). A single-valued, locally bounded, closed-graph map
into a finite-dimensional space is continuous (any subsequential limit of $P(a_n)$ lies in
$m(a)=\{P(a)\}$), and continuity on a compact set is uniform.
\end{proof}

\begin{lemma}[quadratic gain estimate]\label{lem:danskin}
Let $a\notin \X$, $x_a\in m(a)$, $v:=(a-x_a)/d(a)$, $h>0$, and let $y\in m(a+hv)$ be arbitrary.
Then
\[
 d(a+hv)^2\;\ge\;(d(a)+h)^2-\frac{h\,\norm{x_a-y}^2}{d(a)} .
\]
\end{lemma}

\begin{proof}
Since $y$ realizes the distance, $d(a+hv)^2=\norm{a+hv-y}^2=\norm{a-y}^2+2h\sca{v}{a-y}+h^2$.
Now $\norm{a-y}\ge d(a)$, and by the polarization identity
\[
 \sca{a-x_a}{a-y}=\tfrac12\bigl(\norm{a-x_a}^2+\norm{a-y}^2-\norm{x_a-y}^2\bigr)
 \ \ge\ d(a)^2-\tfrac12\norm{x_a-y}^2,
\]
so $\sca{v}{a-y}\ge d(a)-\norm{x_a-y}^2/(2d(a))$. Substituting gives the claim.
\end{proof}

\begin{theorem}\label{thm:fremlinapp}
$\Cx\subseteq\overline{\Mx}$.
\end{theorem}

\begin{proof}
Let $c\in\Cx$, so $R:=d(c)>0$ and $\B(c,R)$ is maximal. Suppose, for contradiction, that
$\overline\B(c,\varepsilon)\cap\Mx=\emptyset$ for some $\varepsilon\in(0,R)$. Write
$\ell:=\varepsilon/2$ and $m_0:=R-\ell>0$. By Lemma~\ref{lem:footmap} the nearest-point map $P$
is single-valued and uniformly continuous on $\overline\B(c,\varepsilon)$; let $\omega$ be a
modulus, i.e.\ $\norm{P(a)-P(a')}\le\omega(h)$ whenever $a,a'\in\overline\B(c,\varepsilon)$,
$\norm{a-a'}\le h$, with $\omega(h)\to0$ as $h\to0$.

\emph{The polygon.} Fix $N\in\mathbb N$, $h:=\ell/N$, and define $a_0:=c$ and
\[
 a_{k+1}:=a_k+h\,v_k,\qquad v_k:=\frac{a_k-P(a_k)}{d(a_k)}\qquad(k=0,\dots,N-1).
\]
By induction $\norm{a_k-c}\le kh\le\ell$, so all nodes stay in $\overline\B(c,\ell)$ and
$d(a_k)\ge R-\ell=m_0$. By Lemma~\ref{lem:danskin} applied at $a_k$ (with
$y=P(a_{k+1})$, $\norm{P(a_k)-P(a_{k+1})}\le\omega(h)$),
\[
 d(a_{k+1})^2\ \ge\ \bigl(d(a_k)+h\bigr)^2-\frac{h\,\omega(h)^2}{m_0} .
\]
Using $\sqrt{A^2-B}\ge A-B/A$ (valid for $0\le B\le A^2$) with $A=d(a_k)+h\ge m_0$:
\[
 d(a_{k+1})\ \ge\ d(a_k)+h-\frac{h\,\omega(h)^2}{m_0\,(d(a_k)+h)}
 \ \ge\ d(a_k)+h\Bigl(1-\frac{\omega(h)^2}{m_0^2}\Bigr),
\]
provided $h$ is small enough that $\omega(h)^2\le m_0^2$. Summing over the $N$ steps,
\[
 d(a_N)\ \ge\ R+\ell\,\bigl(1-\delta(h)\bigr),\qquad \delta(h):=\frac{\omega(h)^2}{m_0^2}
 \xrightarrow[h\to0]{}0 .
\]

\emph{The limit ball.} The endpoints $a_N=a_N^{(h)}$ lie in the compact set
$\overline\B(c,\ell)$; choosing $h_j\to0$ and passing to a convergent subsequence we obtain
$w\in\overline\B(c,\ell)$ with (by continuity of $d$)
\[
 \norm{w-c}\le\ell,\qquad d(w)\ \ge\ R+\ell .
\]
Then $\B(c,R)\subseteq\B(w,d(w))$: for $z\in\B(c,R)$,
$\norm{z-w}\le\norm{z-c}+\ell<R+\ell\le d(w)$. Since $\B(w,d(w))\subseteq E\setminus\X$,
maximality of $\B(c,R)$ forces $(w,d(w))=(c,R)$, contradicting $d(w)\ge R+\ell>R$.
\end{proof}

\begin{remark}
The proof uses neither Brouwer's fixed point theorem nor existence theory for ODEs (Peano), the
two standard tools here: the flow ``always move away from your nearest point'' is replaced by
its Euler polygon, whose per-step loss is controlled by Lemma~\ref{lem:danskin}, and the
inexact gain rate $1-\delta(h)$ is repaired in the limit by compactness. Together with
$\Mx\subseteq\Cx$ (Lemma~\ref{lem:maxballs}(c)) this gives \eqref{eq:fremlin}, and, as noted in
Remark~\ref{rem:motzkin}, Motzkin's theorem on Chebyshev sets follows as a corollary.
\end{remark}

\begin{thebibliography}{9}

\bibitem{Bia} A.~Bia\l{}o\.zyt, \emph{Sets reconstructible with medial axis}, preprint
(corrected version).

\bibitem{Fre} D.~Fremlin, \emph{Skeletons and central sets}, Proc.\ London Math.\ Soc.\
\textbf{74} (1997), no.~3, 701--720.

\bibitem{mathlib} The mathlib Community, \emph{The Lean mathematical library},
\url{https://github.com/leanprover-community/mathlib4}.

\end{thebibliography}

\end{document}
